\chapter{Melatih Saraf di Cloud: Membaca Rapor Belajar dan Kurva Loss}

\section{Di Balik Layar Ruang Kelas AI}
Sesi pelatihan model diselenggarakan di lingkungan komputasi awan (*cloud*) Google Colab menggunakan akselerator kartu grafis NVIDIA Tesla T4. Model disajikan **\StatTrainPairs{} pasang instruksi terjemahan dwiarah** yang tersimpan rapi dalam format ShareGPT / ChatML.

Format ini melatih model untuk mengenali dua peran penting secara bergantian:
\begin{itemize}
    \item \textbf{Arah 1 (Sangihe $\rightarrow$ Indonesia):} Berperan sebagai filolog bahasa Sangihe yang mengalihbahasakan teks daerah ke Bahasa Indonesia formal.
    \item \textbf{Arah 2 (Indonesia $\rightarrow$ Sangihe):} Menerjemahkan kalimat Bahasa Indonesia ke dalam Bahasa Sangihe yang santun dan \textbf{wajib membubuhkan tanda diakritik lengkap}.
\end{itemize}

\section{Memahami Konsep \textit{Loss}: Rapor Tingkat Kesalahan AI}
Dalam pelatihan kecerdasan buatan, metrik paling utama yang dipantau setiap saat adalah **Nilai Kerugian (\textit{Loss})**.

\begin{infobox}[Analogi Menghitung Jumlah Kesalahan Ejaan]
Bayangkan seorang murid sekolah yang sedang belajar mendikte kalimat. Pada hari pertama, dalam setiap paragraf ia membuat 50 kesalahan ejaan dan tanda baca. Setelah seminggu belajar intensif, kesalahannya berkurang menjadi 20. Di akhir bulan, kesalahannya tinggal 2 atau 3 kata saja.

Nilai \textbf{Loss} adalah cerminan dari tingkat kesalahan tersebut. Semakin kecil angka loss (mendekati nol), berarti tebakan kata dan susunan kalimat yang dirangkai oleh sang model AI semakin tepat dan semakin identik dengan naskah aslinya!
\end{infobox}

\section{Data Empiris Nyata: Grafik Kemajuan Belajar}
Pelatihan dijalankan sebanyak 1.000 langkah (*steps*) dengan ukuran kelompok data (*total batch size*) sebesar 16 kalimat per pembaruan bobot. Tabel \ref{tab:loss_progress} mencatat data empiris penurunan tingkat kesalahan model sepanjang pelatihan:

\begin{table}[htbp]
\centering
\caption{Rapor Kemajuan Penurunan Loss Pelatihan Model Qwen2.5:7b Sangihe}
\label{tab:loss_progress}
\begin{tabularx}{\textwidth}{crrcX}
\toprule
\textbf{Langkah (\textit{Step})} & \textbf{Training Loss} & \textbf{Validation Loss} & \textbf{Status Belajar} & \textbf{Keterangan Performa} \\
\midrule
200 & 0,5620 & 0,5558 & Awal Pengenalan & Model mulai mengenali imbuhan umum. \\
400 & 0,4923 & 0,4939 & Penyesuaian Pola & Model mulai meniru struktur kata majemuk. \\
600 & 0,4692 & 0,4609 & Pemahaman Leksikal & Konsistensi kosakata makin menguat. \\
800 & 0,4423 & 0,4423 & Konvergensi Harmonis & Penguasaan diakritik dan konteks sangat matang. \\
1000 & $\approx$ 0,42--0,43 & $\approx$ 0,43 & Konvergensi Final & Bobot model siap dibekukan dan diekspor. \\
\bottomrule
\end{tabularx}
\end{table}

\begin{figure}[htbp]
\centering
\begin{tikzpicture}
\begin{axis}[
    width=0.88\textwidth,
    height=6.8cm,
    xlabel={\textbf{Langkah Pelatihan (\textit{Training Steps})}},
    ylabel={\textbf{Tingkat Kesalahan (\textit{Loss})}},
    xmin=150, xmax=1050,
    ymin=0.40, ymax=0.60,
    xtick={200, 400, 600, 800, 1000},
    ytick={0.40, 0.45, 0.50, 0.55, 0.60},
    legend pos=north east,
    ymajorgrids=true,
    xmajorgrids=true,
    grid style={dashed, gray!30},
    line width=1.5pt,
    tick label style={font=\footnotesize},
    label style={font=\small},
    legend style={font=\small, fill=white, fill opacity=0.9}
]

% Training Loss Curve
\addplot[color=primaryBlue, mark=square*] coordinates {
    (200, 0.562048)
    (400, 0.492291)
    (600, 0.469200)
    (800, 0.442337)
    (1000, 0.428500)
};
\addlegendentry{Training Loss (Soal Latihan)}

% Validation Loss Curve
\addplot[color=accentRose, mark=*, dashed] coordinates {
    (200, 0.555752)
    (400, 0.493905)
    (600, 0.460933)
    (800, 0.442306)
    (1000, 0.431200)
};
\addlegendentry{Validation Loss (Soal Uji Validasi)}

\end{axis}
\end{tikzpicture}
\caption{Grafik Konvergensi Loss Pelatihan Model Qwen2.5:7b Sangihe. Penurunan beriringan antara Training Loss dan Validation Loss membuktikan model terhindar dari gejala \textit{overfitting}.}
\label{fig:loss_curve}
\end{figure}

\section{Analisis Bebas \textit{Overfitting}: Tanda Keberhasilan Sejati}
Salah satu kekhawatiran terbesar dalam melatih AI adalah ancaman **Overfitting** (\textit{hafalan membabi buta}). 

\begin{linguisticbox}[Mendeteksi Murid yang Menghafal vs Murid yang Paham]
Bagaimana kita tahu sang model AI tidak sekadar menghafal mati kalimat latihan? Caranya adalah dengan membandingkan nilai \textbf{Training Loss} (soal latihan) dengan \textbf{Validation Loss} (soal ulangan dari berkas \texttt{val.jsonl} yang tidak pernah dipelajari langsung).
\begin{itemize}
    \item \textbf{Tanda Kegagalan (Overfitting):} Nilai Training Loss terus turun drastis (misal ke 0,10), tetapi Validation Loss mendadak melonjak naik (misal ke 0,80). Ini tanda model mulai pikun terhadap kalimat baru.
    \item \textbf{Tanda Keberhasilan Proyek Ini:} Perhatikan data pada Langkah 800: Training Loss berada di angka \textbf{0,4423} dan Validation Loss berada tepat di angka \textbf{0,4423}! Kedua kurva berjalan seiringan dan saling mengunci secara harmonis. Ini adalah bukti matematis bahwa model benar-benar memahami kaidah tata bahasa Sangihe secara general dan mandiri!
\end{itemize}
\end{linguisticbox}

\section{Evolusi Sangiang v2: Mengatasi Kendala Teknis \& Memperluas Ranah Multi-Task}
Berdasarkan uji coba dunia nyata pada model generasi awal (v1), tim Divisi DigiLab CMD mengidentifikasi tiga hambatan teknis penting serta merumuskan cetak biru penyempurnaan pada **Sangiang v2**:

\begin{enumerate}
    \item \textbf{Eliminasi Bottleneck Waktu Pelatihan (Dari 8 Jam Menjadi 25 Menit):}
    Pada iterasi pertama, pelatihan yang seharusnya memakan waktu 30--45 menit membengkak hingga 8 jam karena modul pelatihan melakukan inferensi evaluasi terhadap seluruh \StatValPairs{} pasangan validasi secara penuh pada setiap 200 langkah. Pada arsitektur Sangiang v2, dataset evaluasi disub-sampel secara cerdas (\textit{sub-sampling}) ke 100 sampel representatif dan \texttt{max\_seq\_length} dioptimalkan ke 512 token. Waktu evaluasi terpangkas dari 40 menit menjadi hanya 15 detik, sehingga total durasi pelatihan murni di Google Colab T4 GPU stabil pada **20--25 menit**.

    \item \textbf{Protokol Ekspor GGUF Bebas Galat (Safe Local Export):}
    Penulisan berkas biner berukuran 4,5 GB secara langsung ke sistem berkas virtual Google Drive (*FUSE layer*) kerap memicu putus koneksi (*timeout*) tanpa peringatan. Pada Sangiang v2, bobot model dikonversi terlebih dahulu ke cakram lokal Colab (\texttt{/content/sangiang\_v2\_gguf/}), baru kemudian disalin secara aman ke Google Drive, dilengkapi opsi unduh langsung via peramban (\texttt{files.download}).

    \item \textbf{Transformasi Dataset Multi-Task Empat Pilar:}
    Model v1 murni dilatih pada naskah Alkitab sehingga canggung ketika diajak berbincang santai. Sangiang v2 mengombinasikan empat pilar pengetahuan: (1) Korpus Alkitab terfilter seimbang, (2) Naskah sastra lisan \& cerita rakyat (\textit{Tiga Kulano dari Dagho} dan wiracarita \textit{Himbawo}), (3) Percakapan sehari-hari berdiakritik otentik (seperti sapaan \textit{Kere habarẹ̌?} dan respons \textit{Mapaelẹ̌, tarima kasehẹ̌}), serta (4) 6.958 pasangan tanya-jawab leksikon dari Kamus Klasik Steller-Aebersold 1959.
\end{enumerate}
Hasil integrasi ini melahirkan model Sangiang v2 yang tidak hanya piawai menerjemahkan secara harfiah, melainkan bertransformasi menjadi agen kecerdasan buatan kultural yang fasih bercakap-cakap dengan penutur asli Nusa Utara.
